\documentclass[12pt,letterpaper,boxed]{hmcpset}
\usepackage[margin=1in,headheight=14pt]{geometry}
\usepackage{amsfonts, amsmath, amssymb, enumerate, fancyhdr, gensymb, lastpage, mathtools, parskip, graphicx}
\usepackage{xcolor, tikz-cd}
\newcommand{\wg}[1]{\textcolor{violet}{#1}}
\newcommand{\OO}{\mathcal O}
\newcommand{\Q}{\mathbb Q}
\newcommand{\R}{\mathbb R}
\newcommand{\C}{\mathbb C}
\newcommand{\Z}{\mathbb Z}
\newcommand{\abs}[1]{\left|#1\right|}
\newcommand{\im}{\text{im }}
\newcommand{\inv}{^{-1}}
\newcommand{\normal}{\unlhd} %% one can also use \trianglelelefteq
\newcommand{\anglee}[1]{\langle #1 \rangle}
\usepackage[shortlabels]{enumitem}

% Numbering macros
\pagestyle{fancy}
\lhead{Will Gilroy}
\chead{Algs Homework \#}
\rhead{03 November 2021}
\lfoot{}
\cfoot{}
\rfoot{Page\ \thepage\ of\ \pageref{LastPage}}

\linespread{1.5}

\newcommand\blankpage{
    \thispagestyle{empty}
    \addtocounter{page}{-1}
    \newpage}
\renewcommand\footrulewidth{0.4pt}

\begin{document}

\problemlist{Algorithms HW } 

%------------------------- Problem 1 -----------------------

\begin{problem}
	\includegraphics[scale=0.8]{1.png}
	\hfill
\end{problem}

\begin{solution}
Since we have uniform convergence on compact subsets we have that $f$
is holomorphic\footnote{
Consider applying Morera's theorem to small
rectangles about each point, and since the convergence is uniform on
compact subsets we can exchange the limit and the integral.
}
on $U$. By isolation of zeroes there's some neighbourhood of $w$
contained in $U$ where $f$ is nonvanishing. Now consider a slightly
smaller neighbourhood about $w$ whose closure is contained in the ball
of nonvanishing $f$. Call this final ball $B$. In particular on
$\partial B$ we have $\abs{f(z)} \geq \delta$ for some positive
$\delta$. On the other hand, for sufficiently large $n$ we have that 
$\abs{f_n(z) - f(z)} < \delta$ for all $z \in \partial B$ since the
convergence is uniform, and so $\abs{f(z)} > \abs{f_z(z) -f(z)}$ on
$\partial B$. We have met the conditions for Morera's theorem and so
it follows that $``f + g'' = f_n(z)$ has the same number of zeroes,
counted with multiplicity, in $B$ as $f(z)$.

By construction, we have that $f$ has exactly one zero in $B$. And so
for each $n$ sufficiently large we have that $f_n$ has exactly one
zero contained in $B$. Call this sequence $\tilde z_n$. 
Since $z_n \in \overline B$ a compact set, we have that some
subsequence of $\tilde z_n$
converges to some point in $z_* \in \overline B$, call this subsequence $z_n$. 
Now, since $f_n \to f$ and $f$ has exactly one zero at $w$ we must
have that $z_n \to w$. 
Indeed 
\[
	0 = \lim_{n \to \infty} f_n(z_n) = f(z_*),
\]
but $f = 0$ only at $w \in \overline B$, and so $z_* = w$. 
The $z_n$ is our desired sequence. 

The argument above only used constructions that were local about the
point $w$ where $f$ vanished, and so I think this conclusion still
holds if $U$ is a non-connected open set, given that $f$ is not
identically zero on any component. 
The existence of the nonvanishing neighbourhood about $w$ depends on
the fact that $f$ is nonvanishing. Given the wording of the
hypothesis, we could have an $f$ which is nonvanishing on all of $U$
but vanishes identically on some connected component. On that
component this argument would not work.


\end{solution}

\newpage

%------------------------- Problem 2 -----------------------

\begin{problem}
	\includegraphics[scale=0.8]{2.png}
	\hfill
\end{problem}

\begin{solution}
Consider for a moment holomorphic functions on the set $X := \C
\setminus \{0,1\}$. 
From our classification of singularities such
functions can either $(1)$ extend to a holomorphic function at either
$0,1$ or both or $(2)$ we can have a pole at either $0$, or $1$. 
\wg{why not both?}. We cannot have an essential singularity at either
$0,1$ because then otherwise, by Castorati-Weirstrass, our function
would not be injective (the image would hit almost every value of $\C$
infinitely often). 

Let us first understand the functions of $f \in Aut(X)$ which extend
to a holomorphic function at both $0,1$. In other words, $f$ is
actually also an automorphism of $\mathbb C$. \wg{wait this might not
be true}

...

Consider for a moment holomorphic functions on the set $X := \C
\setminus \{0,1\}$. Notice that any $f \in Aut(X)$ cannot have an
essential singularity at either $0,1$, otherwise Castorati-Weirstrass
implies that $f$ would not be injective (its image would hit almost
every value of $\C$ infinitely often). That is, any $f \in Aut(X)$
``at worst'' has a pole at $0,1$. It could also extend to a
holomorphic function at either of these points. In other words, a
priori $f: X \to X$ bijective and holomorphic however, we can instead
think about it as a function $f: \hat \C \to \hat \C$.
Moreover, by definition any inverse $f\inv$ is also holomorphic on
$X$. And so by the same argument we can also think of $f\inv$ as a
function $\hat C \to \hat \C$.

First consider the case where $f$ extends to a holomorphic function at
either $0,1$. Since $f$ still needs to be bijective on $X$, if $f$ is
extends to a holomorphic function at $0,1$, then it can only either
fix these points or send them to each other. More precisely, suppose
$f$ extends to a holomorphic function at both $0,1$. Then we must
either have that $f$ fixes both $0,1$ or $f$ exchanges $0,1$. 
The first case is satisfied by only the identity $f(z) = z$ and the
second is satisfied by $f(z) = -z+1$. The latter was found by
recalling that in such a case $f$ is a bijective holomphism of the
entire plane, we found these to be of the form $f(z) = az +b$. 

Now, suppose that $f$ does not extend to a holomorphic function at
both $0,1$. Then we argued above that we can think of $f,f\inv$ as
being extended to $\hat \C \to \hat \C$. Notice that the extended
$f,f\inv$ must still be inverses of each other as functions $\hat \C
\to \hat \C$, since they agree with their original versions on the
dense open set $\hat \C \setminus \{0,1,\infty\}$. 
In other words, we can moreover view $f, f\inv$ as bijections of $\hat
C$. In other words, we can view $f,f\inv$ as Mobius transformations. 

A pole at either $0,1$ corresponds to $f$ sending either $0,1$ to
$\infty \in \hat \C$. But since we argued that $f, f\inv$ must be
bijections on $\hat \C$, it must be that $f$ only has a pole at
exactly one of $0,1$. That is, we can view $f$ as a Mobius
transformation which sends the triple $(0,1,\infty)$ to some
permutation of the same triple. It must send this triple to some
permutation of itself, because we need $f,f\inv$ to be functions on
$X = \C \setminus \{0,1,\infty\}$. 

\wg{now, we should just write down what these can possibly be.}

\end{solution}

\newpage

%------------------------- Problem 3 -----------------------

\begin{problem}
	\includegraphics[scale=0.7]{3.png}
	\hfill
\end{problem}
\begin{solution}
Let us first discuss the boudnary cases. 
When $\abs{z}$ is large we have that this integral goes like
\[
	f(z) \sim \int \zeta^{-(\beta_1 + \beta_2)} d\zeta
\]
In the case where $\beta_1 + \beta_2 \leq 
1$ we have that this integral does not converge and so ``the image of $\infty$
'' is not finite. In other words we get the the image of $\partial \mathbb H$ is
a piecewise union of straight lines, not bounded. In particular, if $\beta_1 +
\beta_2 = 1$ then we have that the angles are both straight lines and we recover
the case of the unbounded strip that we saw in class. On the other hand, if
$\beta_1 + \beta_2 < 1$ then we have that $\beta_1, \beta_2 < 1$ and then we
have an ``obtuse'' unbounded strip.

Now if $\beta_1 + \beta_2 > 1$ we have that $f(\infty)$ is some finite point.
Recall also that $f(\infty)$ and $f(-\infty)$ converge to the same point (we
showed this in class by considering a contour integral of an expanding
semicircle in the upper half plane). 
Consider $\beta_1 + \beta_2 = 2$


\end{solution}

\newpage

%------------------------- Problem 4 -----------------------

\begin{problem}[4]
	\hfill
\end{problem}

\begin{solution}
\end{solution}

\newpage

%------------------------- Problem 5 -----------------------

\begin{problem}
	\hfill
\end{problem}

\begin{solution}
\end{solution}

\newpage

%------------------------- Problem 6 -----------------------

\begin{problem}[4]
	\hfill
\end{problem}

\begin{solution}
\end{solution}

\newpage




\end{document}
